\documentclass[aspectratio=169,12pt]{beamer}
%	\usetheme{Antibes}
	\usetheme{Frankfurt}

%\usepackage{amsmath}

\usepackage{graphics}
\graphicspath{{../img}}

%\definecolor{VerdeApostila}{rgb}{0.058823529411764705, 0.6862745098039216, 0.30980392156862746}
%\usecolortheme[named=VerdeApostila]{structure}

\newcommand{\modo}[2]{
%	#2
} % claro ou escuro
\newcommand{\quebra}{
	\break
}

\newcommand{\iniciar}[3]{

	\begin{frame}[plain]
		\maketitle
		\centering
		
		\textbf{Aula #1}: 
		#2 
		 
		#3
	  \end{frame}

%	\section*{Nesta aula}
	\begin{frame}{Nesta aula}
		\begin{multicols}{3}
			\tableofcontents
		\end{multicols}	
	  \end{frame}  

}

\newcommand{\terminar}[1]{

\section{Exercícios #1}
\begin{frame}[allowframebreaks]{Exercícios}	% containsverbatim
	\input{aula#1_exercícios}	
  \end{frame}	
  
\section*{}  

\begin{frame}[plain]
	\maketitle
	
	\centering	
	\Large
	\textbf{Dúvidas?}
		
	\url{akira.demenech@uel.br}
  \end{frame}

\end{document}

}

%\usepackage[brazil]{babel}	
%\usepackage[T1]{fontenc}

\usepackage{xcolor}

\usepackage{amsmath}
\usepackage{mathdots}

\usepackage{multicol}

\usepackage{hyperref}
\newcommand{\link}[2]{\href{#1}{\texttt{#2}}}

\title{Minicurso de Python}
\author{Guilherme Akira Demenech Mori}
\date{}

\institute{
%	\link{https://www.facebook.com/codificajr/}{
%		\includegraphics[width=0.3\textwidth]{\modo{Logo pronta.pdf}{logo clara.pdf}}
%	}
}

\begin{document}
	\modo{}{\color{white}}	

\iniciar{3}{Funções e Pacotes}{}		  

\subsection{Todos os nomes são Variáveis}

\section{Funções}

	\begin{frame}{O que são funções?}
		Trechos de código armazenados e nomeados.
		
		São executados quando são chamados pelo nome usando \texttt{()}
		
%		\vfill \ \vfill
		\vfill
		
		A depender da definição da função:
		
%		\vfill 
		
%		\hfill 
		pode ser possível inserir parâmetros na função (variáveis de entrada) nos parênteses da chamada.
		
%		\vfill 
		
		\hfill
		e que valores sejam retornados pela função. 
		
%		\vfill				
	\end{frame}
	
	\begin{frame}
%		conteúdo...
		
		
		\textbf{Vantagens}:
		
%		\begin{multicols}{2}
			\begin{itemize}
				\item Redução da repetição  
				\item Fácil alteração 
				\item Simplificação 
			\end{itemize}
%		\end{multicols}
		
	\end{frame}

	\subsection{Funções nativas}
	\begin{frame}
%		conteúdo...	
	\begin{multicols}{3}
%		conteúdo...
	
		\begin{itemize}
			\item \texttt{print} 
			\item \texttt{input}	
			
		\vfill 
			
			\item \texttt{str}	
			\item \texttt{bool} 
			\item \texttt{int}, \texttt{float} e \texttt{complex} 
			\item \texttt{list} e \texttt{tuple}
			\item \texttt{dict} e \texttt{set}
	
		\vfill
	
			\item \texttt{hex}
			\item \texttt{oct}
			\item \texttt{bin}
		
		\vfill 	
		
			\item \texttt{len}		
			\item \texttt{abs}
			\item \texttt{round}

		\vfill 
	
			\item \texttt{id}

				
			\item \texttt{type}
			\item \texttt{dir}
			
		\vfill	
			
			\item ....
	
		\end{itemize}
	\end{multicols}
	\end{frame}
	
	\subsection{Definição}
	
	\begin{frame}[containsverbatim]{Definição}
%		conteúdo...

		A definição da função deverá ter um nome válido como variável.
		
		As ações podem incluir estruturas condicionais, laços de repetição, chamadas de funções e quaisquer operações. 

		\begin{verbatim}
			def nome_da_funcao ():
				
				acoes
		\end{verbatim}

	\end{frame}
	
	\begin{frame}[containsverbatim]
		%		conteúdo...
		
		Utilize \texttt{global} para acessar e alterar variáveis criadas fora da função.
		
		\begin{verbatim}
			def funcao ():
			
				global v
				
				acoes				
			
		\end{verbatim}
		
	\end{frame}
	
	\begin{frame}[containsverbatim]
		%		conteúdo...
		
		Utilize \texttt{return} para retornar algum valor para o local de chamada.				
		
		As variáveis entre parênteses na declaração são recebidas como parâmetros.
		
		\begin{verbatim}
			def nome_da_outra (x, y, z):
			
				acoes
				
				return valor
		\end{verbatim}
		
	\end{frame}
	
	
	
	\subsection{Uso}
	
	\begin{frame}[containsverbatim]{Uso}
		%		conteúdo...
		
		Para chamar uma função, use o nome dela e coloque parênteses depois.
		
		\begin{verbatim}
			nome_da_funcao()						
			
			w = nome_da_outra(a,b,c)
		\end{verbatim}
		
	\end{frame}
	
	\subsection{Funções também são Variáveis}

\section{Pacotes}

	\subsection{Pacotes nativos}
	
	\begin{frame}
		%		conteúdo...	
		\begin{multicols}{3}
			%		conteúdo...
			
			\begin{itemize}
				\item \texttt{time}
				\item \texttt{tkinter}
				
				\vfill 
				 
				\item \texttt{sys}									 				
				\item \texttt{os}	
				
				\vfill 
				
				 
				\item \texttt{math} 					
				
				\item ....
				
			\end{itemize}
		\end{multicols}
		
		\vfill \hfill 
	\end{frame}	
	
	\begin{frame}{Importação de pacotes}
		
		Importe o pacote que você quer utilizar no seu programa com 
		
		\texttt{import nomedopacote}
		
		E, depois, para utilizar funções dele, utilize o nome e \texttt{.}
		
		\texttt{nomedopacote.funcao()}
	\end{frame}
	
	\subsection{Instalação de pacotes}
	
	\begin{frame}
		%		conteúdo...	
		\begin{multicols}{3}
			%		conteúdo...
			
			\begin{itemize}
				\item \texttt{numpy}
				\item \texttt{scipy}
				
				\vfill 
				
				\item \texttt{scikit-learn}									 				
				\item \texttt{pandas}	
				
				\vfill 
				
				
				\item \texttt{matplotlib} 					
				
				\item ....
				
			\end{itemize}
		\end{multicols}
		
		\vfill  				
		
		\hfill
		\texttt{python3 -m pip install nomedopacote}
		
		\hfill
		\texttt{python -m pip install nomedopacote}
		
		\hfill
		\texttt{py -m pip install nomedopacote}
		
	\end{frame}	

	\begin{frame}
%	conteúdo...
	
		
		Se o seu sistema recusar a instalação de pacotes via \link{https://packaging.python.org/en/latest/tutorials/installing-packages/}{pip}, você talvez precise usar ambientes virtuais 
		(\link{https://docs.python.org/pt-br/3/library/venv.html}{venv}) 		
		ou instalar diretamente pelo gerenciador do sistema 
		(\texttt{apt}, por exemplo)
		
	\end{frame} 
	
	 
	
	\subsection{Pacotes também são Variáveis}

\terminar{3}
